(1) 均匀旋转流体的速度场为 \(\mathbf u=\Omega(-y,x)\)。一个小惯性颗粒满足 \(\dot{\mathbf r}=\mathbf v,\ \tau\dot{\mathbf v}=\mathbf u(\mathbf r)-\mathbf v\)，其中 \(\Omega,\tau>0\)。写成二维线性系统并推导复坐标 \(z=x+iy\) 的常系数方程。

(2) 求完整的四个实动力学特征根及复模态；判断粒子是否能在旋涡中心渐近稳定，解释Stokes数 \(\mathrm{St}=\Omega\tau\) 的作用。

(3) 求 \(\mathrm{St}\ll1\) 时控制长时间径向迁移的特征根至 \(O(\mathrm{St}^2)\)，推导径向漂移与角速度的首阶表达式。

(4) 对 \(\Omega=0\) 或 \(\tau\to0^+\) 的极限做物理检验；说明此模型为什么不能把颗粒速度在每个瞬间直接设置为背景流速。
